% concatenated from JSL_well_defined_not_predetermined_no_self_citations.tex
\documentclass[11pt]{article}

\usepackage[a4paper,margin=1in]{geometry}
\usepackage{amsmath,amssymb,amsthm,mathtools}
\usepackage{enumitem}
\usepackage[hidelinks]{hyperref}
\usepackage[nameinlink,noabbrev]{cleveref}
\usepackage{microtype}

\newtheorem{definition}{Definition}[section]

\newtheorem{corollary}{Corollary}[definition]

\newtheorem{axiom}[definition]{Axiom}
\newtheorem{proposition}[definition]{Proposition}
\newtheorem{theorem}[definition]{Theorem}
\newtheorem{lemma}[definition]{Lemma}
\newtheorem{remark}[definition]{Remark}
\newtheorem{example}[definition]{Example}
\newtheorem{question}[definition]{Question}

\newcommand{\bits}{\{0,1\}}
\newcommand{\strings}{\bits^*}
\newcommand{\Aev}{\mathcal A}
\newcommand{\World}{\mathcal W}
\newcommand{\Hist}{\mathcal H}
\newcommand{\Obs}{\mathcal O}

\title{\textbf{Well-Defined but Not Predetermined}\\
\large Persistent Mathematical Objects, Embedded Observers, and the Provability of \(P=NP\)}

\author{Rasoul Ramezanian\\ \small University of Fribourg \\ \small rasoul.ramezanian@unifr.ch}
\date{}

\begin{document}
\maketitle

\begin{abstract} 
We develop a foundational distinction between \emph{well-definedness} and
\emph{predetermination}.  A persistent mathematical or computational object may return a
unique answer whenever a query is made and preserve every previously returned answer, while
its answers to as-yet unqueried inputs are not represented by one completed,
history-independent extension.  We call such an object \emph{well-defined but
non-predetermined}.

The framework is developed through persistent, history-indexed partial functions and set
processes.  Every realized history of such a process is extensionally compatible with a
classical predetermined object, although different incompatible histories may require
different classical completions.   

We next formulate an observer-based semantics.  All observers inhabit one shared
computational world, their interactions modify one common global state, and they have access
to one realized history.  Under axioms of shared history, historical irreversibility, persistence, and
historical indistinguishability, embedded observers cannot establish from the realized
history alone whether its compatible extension was predetermined or progressively formed.

A persistently evolutionary polynomial time computable oracle supplies a concrete example and yields a
deterministic--nondeterministic separation under an explicitly history-sensitive
  semantics.  We then derive a conditional result concerning the provability of \(P=NP\). Suppose that an observer cannot distinguish its actual computational world from a \textit{persistently evolutionary world} in which \(P\neq NP\). Suppose further that the observer's entire internal reasoning process—including any purported proof of \(P=NP\)—is reproduced in that evolutionary world. The same proof process would then lead the observer to conclude \(P=NP\) in a world in which \(P=NP\) is false. Consequently, that process cannot constitute a sound proof of \(P=NP\). Under the proposed observer axioms and the existence of such an indistinguishable evolutionary realization, no embedded observer can therefore possess a sound internal proof of \(P=NP\). If human mathematicians are themselves embedded observers in this sense, the same conclusion applies to us.


The conclusion is an observer-relative unprovability theorem, not an unconditional proof of
classical \(P\neq NP\). 
\end{abstract}

\noindent\textbf{Keywords:} non-predetermined functions; persistent computation; embedded
observers; irreversible histories; \(P\) versus \(NP\); provability; simulation,  Brouwer Choice Sequences.

\medskip
%\noindent\textbf{MSC 2020:} 03D10, 03D15, 03B70, 68Q15.

\section{Introduction}

Classical mathematics normally treats a function or set as extensionally complete.  If
\(f:X\to Y\), then every value \(f(x)\) belongs to one already determined completed graph.  If
\(A\subseteq X\), then every membership proposition \(x\in A\) has a fixed truth value,
independently of whether that question is ever asked.

This ordinary conception combines two logically different properties:
\begin{enumerate}[label=(\roman*)]
  \item \emph{well-definedness}: an input has a unique value;
  \item \emph{predetermination}: all values belong to one completed extension fixed
  independently of \textit{interaction} history.
\end{enumerate}

The purpose of this paper is to separate these properties.  A process may determine a unique
answer when an input is first encountered and preserve that answer forever, while the answer
to a previously unencountered input remains sensitive to the order of future interactions.
Such a process is well-defined along every realized history, but no single
history-independent extension need represent all counterfactual histories.

\begin{quote}
 \textit{Well-definedness may arise from the persistence of already determined values rather than from the predetermination of all future values, while an embedded observer is unable to distinguish whether the values are progressively determined and preserved or fixed in advance.
}
\end{quote}

We develop a framework to formulate explicitly the assumptions concerning embedded observers, and to distinguish carefully between the following three claims:
\begin{enumerate}[label=(\arabic*)]
  \item there exist coherent formal models of persistent non-predetermined computation in which deterministic and nondeterministic polynomial-time computation are separated;
  \item embedded observers are unable to distinguish such non-predetermined models from predetermined ones on the basis of their realized history;
  \item under additional assumptions concerning the preservation and soundness of internal proofs, such observers cannot soundly prove
  \(P=NP\).
\end{enumerate}

The third claim is epistemic.  It is not, by itself, the classical mathematical conclusion
\(P\neq NP\).

The framework has affinities with mathematical un-finished objects as Brouwer's choice sequences
\cite{TroelstraVanDalen1988,VanAtten2007}, interactive computation
\cite{Goldin2000,GoldinSmolkaWegner2006,GoldinWegner2008}, Bostrom's simulation arguments
\cite{Bostrom2003}, and   process-based descriptions of evolving natural systems~\cite{WongEtAl2023}.   

 The main theorem formalizes the following argument. Suppose that an observer cannot distinguish its actual world from a persistently evolutionary world in which \(P\neq NP\). Suppose, moreover, that the observer would carry out exactly the same reasoning and obtain exactly the same purported proof of \(P=NP\) in both worlds. Since \(P=NP\) is false in the evolutionary world, this reasoning cannot constitute a sound proof of \(P=NP\). Thus, under the assumptions introduced in this paper, an embedded observer cannot possess a sound proof of \(P=NP\).

 \section{Related Work}
\label{sec:related-work}

\subsection{Brouwer's choice sequences}

The idea that a mathematical object need not be given as a completed totality from the outset
has a well-established precedent in intuitionistic mathematics.  In Brouwer's theory of
choice sequences, an infinite sequence is conceived as being progressively generated rather
than as a completed infinite object whose entries are all fixed in advance; see, for example,
\cite{TroelstraVanDalen1988,VanAtten2007}.  At any finite stage, only an initial segment of
the sequence has been determined, while further entries may be determined as the construction
continues.

This viewpoint is closely related to the distinction between well-definedness and
predetermination considered in the present paper.  A persistent non-predetermined function
similarly develops through time: once a value has been determined, it remains fixed, while
values not yet determined need not belong to a completed extension fixed independently of
future interactions.  Thus, in both settings, mathematical definiteness need not be identified
with the prior existence of a completed infinite object.

There are, however, important differences.  Our framework is not intended as a reformulation
of Brouwerian intuitionism or of the theory of choice sequences.  Here the central object is
an interactive computational process whose future extension may depend on the history and
order of interactions.  Moreover, our main concern is observer-relative: an embedded observer
has access only to one realized history and, under the assumptions introduced in this paper,
may be unable to determine whether the observed values were fixed in advance or progressively
determined and subsequently preserved.  Choice sequences therefore provide an important
foundational precedent for mathematical objects that unfold over time, while the present
framework develops this idea in a persistent, computational, and observer-dependent setting.

\subsection{Simulation and embedded reasoning}

The simulation hypothesis is not assumed in this paper.  Bostrom's simulation argument
concerns probabilistic alternatives involving the development of posthuman civilizations and
the possibility that observers like ourselves inhabit ancestor simulations
\cite{Bostrom2003}.  The present framework requires none of these claims.  It draws only on
the more general methodological possibility that observers may have access to the behavior
generated within a world without having access to the external mechanism by which that world
is realized.

This distinction is particularly relevant when the observer is itself part of the system
under investigation.  In the present framework, observing, reasoning, recording information,
and constructing or verifying a mathematical proof are themselves events occurring within the
modeled world.  The observer therefore does not occupy an external standpoint from which
different realizations of the same history can be inspected and compared.

For this reason, three notions must be kept distinct:
\begin{enumerate}[label=(\roman*)]
\item equality of the histories accessible to the observer;
\item preservation of a particular internal reasoning or proof process;
\item equality of mathematical truth across the underlying worlds.
\end{enumerate}
Two worlds may be indistinguishable from the observer's internal perspective without thereby
being identical as external mathematical structures.  This distinction is central to the
observer-relative argument developed below.

\subsection{Scientific motivation from evolving systems}

Wong et al.\ propose a general account of evolving systems based on combinatorial richness,
the generation of configurations, and selection for function \cite{WongEtAl2023}.  Several
aspects of their framework provide useful conceptual motivation for the present work.

First, their distinction between static and dynamic persistence emphasizes that what persists
through time need not be an unchanging configuration; persistence may instead characterize an
ongoing process.  This is conceptually related to our use of persistence, although persistence
has a more specific meaning here: once a value or membership decision has been realized, it
remains fixed throughout every continuation of the realized history.

Second, Wong et al.\ emphasize the temporal asymmetry and context dependence of evolving
processes.  This is compatible with, although weaker than, the history dependence considered
here, where the future extension of a mathematical object may depend on the order in which
interactions occur.

Third, their discussion of novelty generation and open-ended evolution includes the
possibility that evolving systems construct new possibility spaces rather than merely explore
a fixed space of configurations.  This provides useful motivation for distinguishing a
predetermined but incompletely known future from a process in which future states are
progressively formed.

These connections are motivational rather than deductive.  The theory of evolving systems
does not imply persistent non-predetermination in the technical sense used in this paper, nor
does it establish the observer-indistinguishability assumptions introduced below.

\subsection{Relation to standard complexity barriers}

The relationship between the present construction and classical barriers in complexity theory
requires particular care.  Baker, Gill, and Solovay showed that there exist fixed oracles
relative to which \(P=NP\) and other fixed oracles relative to which \(P\neq NP\)
\cite{BakerGillSolovay1975}.  Consequently, any proof technique that relativizes cannot by
itself resolve the classical \(P\) versus \(NP\) problem.  Natural proofs and algebrization
identify additional limitations on broad classes of proof techniques
\cite{RazborovRudich1997,AaronsonWigderson2009}.

The polynomial-time computable oracle considered in the present paper, which separates deterministic from nondeterministic polynomial-time computation under the proposed semantics, differs fundamentally from a classical fixed oracle.  Its answers are persistent once determined, but its extension is formed through
interaction and may depend on the history and order of queries.  The resulting complexity
classes are therefore defined under an explicitly history-sensitive computational semantics.
The evolutionary separation established in Theorem~\ref{thm:evsep} should consequently not
be interpreted as an ordinary relativized separation for a fixed oracle.

Accordingly, the present framework does not claim that the evolutionary separation alone
constitutes a proof of the classical statement \(P\neq NP\).  Rather, the separation provides
the computational component of the subsequent observer-relative argument.   



\section{Well-definedness without predetermination}

\subsection{Histories and partial functions}

Let \(X\) be a countable set of possible inputs and \(Y\) a set of possible outputs.

\begin{definition}[History]
A finite history is a sequence
\[
  H=((x_1,y_1),\ldots,(x_t,y_t))\in(X\times Y)^{<\omega}.
\]
It is \emph{consistent} if \(x_i=x_j\) implies \(y_i=y_j\).  If \(H'\) extends \(H\),
we write \(H\preceq H'\).
\end{definition}

Every consistent history determines a partial function \(f_H\).

\begin{definition}[Persistence]
A history-generating process is \emph{persistent} if
\[
  H\preceq H'\quad\Longrightarrow\quad f_H\subseteq f_{H'}.
\]
\end{definition}

Thus every value once realized remains fixed.

\begin{definition}[Persistent non-predetermination]
A persistent process is \emph{non-predetermined} if there exist admissible histories
\(H,H'\) and an input \(x\) such that
\[
  x\in\operatorname{dom}(f_H)\cap\operatorname{dom}(f_{H'})
  \quad\text{and}\quad
  f_H(x)\neq f_{H'}(x).
\]
\end{definition}

The disagreement occurs across incompatible histories, never inside one realized history.

\subsection{A canonical example}

\begin{example}
Let \(W\) initially be empty.  When \(n\in\mathbb N\) is presented, define
\[
g(n)=
\begin{cases}
z, & (n,z)\in W,\\
|W|+1, & \text{otherwise},
\end{cases}
\]
and in the second case update
\[
W\leftarrow W\cup\{(n,|W|+1)\}.
\]
If the arguments arrive in the order
\[
7,9,1,11,
\]
then \(g\) grows as follows:
\[
g(7)=1,\quad g(9)=2,\quad g(1)=3,\quad g(11)=4.
\]
If they arrive in the order
\[
9,1,7,11,
\]
then \(g\) forms as follows:
\[
g(9)=1,\quad g(1)=2,\quad g(7)=3,\quad g(11)=4.
\]
Repeated presentation of an argument always returns the value previously assigned to it.
\end{example}

The process is not random in the relevant sense: given the current state and next input, the
next output is fixed.  What is absent is a single extension independent of the future input
history. The function \(g\) is a mathematical object whose values are progressively determined over time, in a manner analogous to Brouwer's choice sequences, rather than being specified by a completed extension from the outset. 

 

\subsection{Persistent set processes}

Let \(\Sigma=\bits\).  At history \(H\), define
\[
A_H^+=\{x\in\strings:\chi_H(x)=1\},\qquad
A_H^-=\{x\in\strings:\chi_H(x)=0\},
\]
and
\[
U_H=\strings\setminus(A_H^+\cup A_H^-).
\]

\begin{definition}[Persistent set process]
A persistent set process is a family of partial characteristic functions
\[
  \chi_H:\operatorname{dom}(\chi_H)\to\bits
\]
indexed by admissible histories, such that
\[
  H\preceq H'\quad\Longrightarrow\quad \chi_H\subseteq\chi_{H'}.
\]
\end{definition}

\begin{proposition}[Finite trace completion]\label{prop:completion}
For every finite consistent history \(H\), there exists a classical set
\(B_H\subseteq\strings\) whose membership answers agree with all observations in \(H\).
\end{proposition}

\begin{proof}
Choose any \(C\subseteq U_H\) and define
\[
  B_H=A_H^+\cup C.
\]
For every queried string \(x\), membership in \(B_H\) agrees with the answer in \(H\).
\end{proof}

\begin{remark}
The proposition establishes extensional compatibility along one history.  It does not provide
one static set valid across all possible histories.
\end{remark}

 

 

 

\section{Observers in one shared computational world}

Let
\[
  \World=(S,s_0,\tau)
\]
be a computational world.  Observers, computers, records, and proof activities are processes
within \(\World\).  A realized history is
\[
s_0\xrightarrow{e_1}s_1\xrightarrow{e_2}s_2\longrightarrow\cdots .
\]

\begin{axiom}[Shared history]\label{ax:shared}
All observers in one world interact with one common global state.  An interaction by one
observer becomes part of the state encountered by every later observer.
\end{axiom}

\begin{axiom}[Historical irreversibility]\label{ax:irreversible}
An embedded observer cannot return the entire world to a previous total state (s) while retaining information acquired after that state. Restoring the entire world to (s) would necessarily restore the observer, including their memory and all available records, to their state at (s). Consequently, the observer cannot experience two different continuations of exactly the same total past state and retain enough information to compare them.
\end{axiom}

 

\begin{axiom}[Persistence of realized facts]\label{ax:persist}
Once the world has realized a membership or functional value, every later realized state
preserves it.
\end{axiom}

\begin{definition}[Internal historical equivalence]
Two world-histories are internally historically equivalent for an observer \(O\) if the
observer has the same accessible sequence of observations, actions, memories, written symbols,
proof states, and verification states in both.
\end{definition}

\begin{axiom}[Historical indistinguishability]\label{ax:ind}
If two worlds generate internally indistinguishable histories for an observer (O), then, on the basis of that history alone, (O) cannot determine which of the two worlds it inhabits.
\end{axiom}

%The axiom does not say that the worlds satisfy the same external formulas.  It says that the observer cannot identify the implementation from the shared internal trace.

\subsection{Why many observers do not defeat indistinguishability}

The observer argument does not depend on there being only one observer.  Let
\[
  \Obs=\{O_1,O_2,\ldots\}
\]
be any population of observers.  By Axiom~\ref{ax:shared}, their interventions form one
interleaved global history.  If \(O_1\) changes the state, \(O_2\) encounters the changed
state.

To distinguish predetermination directly from persistent formation, the observers would need
to compare incompatible continuations from the numerically identical total past while
retaining information from both continuations.  This is what Axiom~\ref{ax:irreversible}
excludes.

%Creating independent copies is not equivalent.  Copies are distinct systems, not two simultaneously remembered futures of one numerically identical total state.



\section{A persistently evolutionary polynomial time  oracle}\label{sec:oracleA}
In this section, we introduce a  persistently evolutionary polynomial time computable oracle denoted by \(\Aev\).

Let
\[
M=(Q,\bits,\delta,q_0,F)
\]
be a partial deterministic finite automaton.  Initially
\[
Q=\{q_0\},\qquad \delta=\varnothing,\qquad F=\varnothing.
\]
where \(Q\) shows states of the automaton, \( \delta\)  denotes the partial transition function, and  \(F\) refers to final states.



On a query \(x=a_1\cdots a_n\), follow the defined transitions.

 \begin{enumerate}[label=\textbf{Rule \arabic*.}]
    \item If, following transitions in \( \delta\), all of \(x\) is read and the final state lies in \(F\), accept without change.
    \item If, following transitions in \( \delta\), all of \(x\) is read and the final state \(q\notin F\), reject if \(q\) has a
    one-symbol transition to an accepting state; otherwise add \(q\) to \(F\) and accept.
    \item If, following transitions in \( \delta\),  the automaton crashes on an undefined transition, add a new path for the unread
    suffix (adding new states and transitions) to complete reading \(x\), and make the final new state accepting.
\end{enumerate}

Call the resulting process \(\Aev\).



\begin{proposition}[Persistence]\label{prop:persistent-oracle}
Every membership answer returned by \(\Aev\) is permanent.
\end{proposition}

\begin{proof}
An accepted string ends at an accepting state, whose status is never removed.  A rejected
string ends at a nonaccepting state having a one-symbol transition to an accepting state;
neither that transition nor the nonaccepting status is later removed.
\end{proof}

\begin{proposition}[History dependence]\label{prop:history}
The extension generated by \(\Aev\) is not predetermined.
\end{proposition}

\begin{proof}
If \(10\) is queried first, Rule 3 creates an accepting path for \(10\).  If \(101\) is
queried first, the state reached after \(10\) is nonaccepting and has a one-symbol transition
to an accepting state, so a later query \(10\) is rejected.
\end{proof}

\begin{proposition}[Response cost]
Under unit-cost access to transitions and unit cost per state or transition update, every query
\(x\) is processed in \(O(|x|)\) operations.
\end{proposition}

\begin{proof}
At most \(|x|\) transitions are traversed. If a path is added, at most \(|x|\) states and
transitions are created.
\end{proof}

The automaton above is the formal representation of \(\mathcal A\) used
throughout the paper.  It is also possible to realize the same process as a
persistently growing program.  In this representation, every finite binary
prefix corresponds to a directory containing a constant-size
\texttt{node.py} file with three binary values recording whether the
corresponding state is accepting and whether its \(0\)- and \(1\)-branches
have already been generated.  Processing an input follows only the
directories corresponding to its successive prefixes, while Rules~2 and~3
modify or extend this code when necessary.  Appendix~\ref{app:growing-code}
describes this implementation and shows that each query is processed in
time polynomial in the length of the input.

\section{History-sensitive complexity}

A deterministic machine interacts sequentially with the current global state of \(\Aev\).
Let \(P_{\mathrm{ev}}^{\Aev}\) denote languages decided in polynomial time equipped with the oracle \(\Aev\).

For nondeterminism, we use certificates.

\begin{definition}
A language \(L\) belongs to \(NP_{\mathrm{ev}}^{\Aev}\) if there are
\(J\in P_{\mathrm{ev}}^{\Aev}\) and a polynomial \(q\) such that
\[
x\in L
\quad\Longleftrightarrow\quad
\exists y\bigl(|y|\le q(|x|)\land (x,y)\in J\bigr).
\]
\end{definition}

Let
\[
L_{\Aev}=\{x:\Aev\text{ accepts }x\}
\]
 and
\[
L_{\exists}
=
\{w:\exists x\,(|x|=|w|\land x\in L_{\Aev})\}.
\]

\begin{proposition}\label{prop:belongtoNP}
\(L_{\exists}\in NP_{\mathrm{ev}}^{\Aev}\).
\end{proposition}

\begin{proof}
Use \(x\) as the certificate, verify \(|x|=|w|\), and query \(\Aev\).
\end{proof}

\begin{definition}[Continuation-robust decision]
A decision made at history \(H\) is continuation-robust if it remains correct in every
admissible extension \(H'\succeq H\).
\end{definition}

\begin{theorem}[Evolutionary separation]\label{thm:evsep}
Under continuation-robust semantics,
\[
P_{\mathrm{ev}}^{\Aev}\neq NP_{\mathrm{ev}}^{\Aev}.
\]
\end{theorem}

\begin{proof}
[Proof sketch]
By Proposition~\ref{prop:belongtoNP}, \(L_{\exists}\in NP^{\Aev}_{\mathrm{ev}}\).  
 Suppose, toward a contradiction, that a deterministic polynomial-time procedure \(D\) continuation-robustly decides \(L_{\exists}\). Let \(p\) be a polynomial bounding the running time of \(D\). Choose \(n\) sufficiently large so that
 \( p(n)<2^n \) 
 and no string of length at least \(n\) has previously been queried to \(\Aev\) by any embedded observer. 
 
 Fix an arbitrary \[ w\in\{0,1\}^n. \] We run \(D\) on \(w\) in the shared computational world. All observers and computational procedures interact with the same global state of \(\Aev\). Hence, while \(D\) is running, other embedded observers may interleave their interactions with those of \(D\) and submit additional queries to \(\Aev\). These queries become part of the same persistent global history. 
 
 Whenever \(D\) is about to submit an oracle query \(y\) with \(|y|\ge n\), let \(x\) be the prefix of \(y\) of length \(n\). If no previous interaction has already determined the relevant branch through \(x\), another embedded observer first queries \(x0\). By the rules defining \(\Aev\), this makes \(x0\) accepted and causes \(x\) itself to be rejected permanently. The query of \(D\) is then processed in the resulting global state. Since \(D\) terminates within \(p(n)\) steps, it makes at most \(p(n)\) oracle queries. Consequently, the above construction can involve at most \(p(n)\) distinct length-\(n\) prefixes. Because \( p(n)<2^n, \) at least one string \( z\in\{0,1\}^n \) remains untouched: no query occurring during the computation has \(z\) as its length-\(n\) prefix. After \(D\) terminates, there are two cases.
 
 If \(D\) rejects \(w\), choose such an untouched string \(z\). Since no previous query has passed through the branch represented by \(z\), querying \(z\) next causes Rule~3 to create the missing path and make \(z\) accepting. Hence \( z\in L_{\Aev}. \) Since \(|z|=|w|=n\), it follows that \( w\in L_{\exists}. \) Thus there is an admissible continuation of the shared global history in which the negative answer of \(D\) is false, contradicting continuation robustness. 
 
 If \(D\) accepts \(w\), then every length-\(n\) string that arose through the above interleaving has already been made permanently rejected. For each remaining untouched string \( x\in\{0,1\}^n, \) another embedded observer may query \(x0\) before \(x\). By the same rule, this permanently forces \(x\) to be rejected. After these finitely many   interactions, every string of length \(n\) is rejected, so \[ L_{\Aev}\cap\{0,1\}^n=\varnothing. \] Therefore no witness of length \(n\) exists and hence \( w\notin L_{\exists}. \) 
 Thus there is an admissible continuation in which the positive answer of \(D\) is false, again contradicting continuation robustness.
 
 The essential point is that a polynomial-time procedure can interact with only polynomially many parts of an exponentially large space of possible length-\(n\) witnesses. When \(D\) terminates, it therefore cannot have fixed enough of the future evolution of the shared global state to make either answer invariant under all admissible continuations. Hence no deterministic polynomial-time procedure continuation-robustly decides \(L_{\exists}\). 
 
 Since \( L_{\exists}\in NP^{\Aev}_{\mathrm{ev}}, \) we conclude that \[ P^{\Aev}_{\mathrm{ev}} \neq NP^{\Aev}_{\mathrm{ev}}. \]
The complete argument is given in Appendix~\ref{app:evolutionary-separation}.
\end{proof}

 
The proof rests on two related ideas. First, the global state is shared:
\(D\) does not compute in isolation, and a proposition established by an
embedded observer must remain correct regardless of how other embedded
observers may simultaneously interleave their interactions with the same global
state. Second,   a polynomial-time procedure can
interact with only polynomially many parts of an exponentially large space
of possible length-\(n\) witnesses. When \(D\) terminates, it therefore
cannot have fixed enough of the future evolution of the shared global state
to make either answer invariant under all admissible continuations. In this
sense, \(D\) cannot close the relevant future. Hence no deterministic
polynomial-time procedure continuation-robustly decides \(L_{\exists}\).
 


\begin{remark}
This theorem concerns the explicitly defined evolutionary semantics.  It is not a standard
relativized separation for a fixed oracle set.  The distinction is essential in view of the
relativization barrier \cite{BakerGillSolovay1975}.
\end{remark}

\section{Proof as an internal event}

A formal derivation available to an embedded observer must be written, stored, read, checked,
or communicated.  These activities are events in the world-history.

\begin{definition}[Internal proof realization]
An internal realization of a finite derivation \(\pi\) consists of the observer's accessible
sequence of symbol states, rule applications, memory states, and verification states through
which \(\pi\) is produced and accepted as a proof.
\end{definition}

\begin{axiom}[Internal proof preservation]\label{ax:proof}
Suppose an observer has internally historically equivalent histories in two worlds, and the
same finite derivation \(\pi\), formal axioms, inference rules, intended semantics, and
verification process are reproduced in both.  If \(\pi\) is sound under that intended
semantics, its conclusion cannot be false in one of the two worlds.
\end{axiom}

%This axiom does not say that historically indistinguishable worlds satisfy every same sentence. It concerns a particular proof process whose complete internally relevant realization is preserved.

\subsection{Observer-compatible evolutionary realizations}

\begin{definition}[Observer-compatible realization]
Let \(O\) be an embedded observer with realized history \(H\). A world
\(\World'\) is an \emph{observer-compatible realization} of \(H\) if
\(O\)'s complete internally accessible history in \(\World'\) is identical
to \(H\). In particular, \(O\) has the same observations, memories,
interactions, records, and internally realized reasoning processes in both
worlds.
\end{definition}


\begin{proposition}[Finite-history realization]
\label{prop:finite-realization}
Let \(H\) be any finite consistent history accessible to an embedded
observer \(O\). Then there exist both a predetermined realization and a
persistent non-predetermined realization that are observer-compatible with
\(H\).
\end{proposition}

\begin{proof}
Since \(H\) is finite and consistent, the values determined along \(H\)
define a finite partial function. A predetermined realization is obtained
by extending this partial function arbitrarily to a total function on the
entire domain.

For the non-predetermined realization, preserve every value already fixed
by \(H\), while leaving values not determined by \(H\) to be assigned by a
persistent evolutionary process. The process agrees with \(H\) on every
interaction already observed by \(O\), but previously undetermined values
may be assigned differently in different admissible continuations.
Therefore the two realizations generate the same internally accessible
history for \(O\), although one is predetermined and the other is not.
\end{proof}


%Condition (iv) is substantive.  It cannot be replaced merely by the observation that \(P_{\mathrm{ev}}^{\Aev}\neq NP_{\mathrm{ev}}^{\Aev}\) if the purported proof concerns different, classical complexity classes.  A fully formal application must align the sentence proved with the sentence falsified. This alignment requirement identifies the central technical burden of the argument.

%\section{The main observer-relative theorem}

\begin{theorem}[Embedded-observer unprovability]\label{thm:main}
Assume axioms~\ref{ax:shared},\ref{ax:irreversible},\ref{ax:persist},\ref{ax:ind},\ref{ax:proof}.  Let \(O\) be an embedded
mathematical observer.  Suppose \(O\)'s realized history has an observer-compatible
evolutionary realization in the sense above.  Then \(O\) cannot possess a sound internal proof
of \(P=NP\).
\end{theorem}

\begin{proof}
Assume that \(O\) possesses a sound internal proof \(\pi\) of \(P=NP\). 

By Theorem~\ref{thm:evsep},
there is an observer-compatible evolutionary realization in which \(O\)'s complete internal
history and the proof realization of \(\pi\) are reproduced, with the same formal system,
intended semantics, and verification process.

By Axiom \ref{ax:proof}, a sound internal proof
process reproduced with the same intended semantics cannot establish a sentence false in one
of the two worlds.  Therefore \(\pi\) is not sound, contrary to assumption.
\end{proof}



\subsection{What follows, and what does not}

The result is conditional.
The theorem has the logical form
\[
\mathsf{Axioms}
\quad\Longrightarrow\quad
\text{no sound internal proof of }P=NP.
\]
The antecedent is not established by ordinary complexity theory.

 Unprovability does not imply falsity.
In general,
\[
\nvdash P=NP
\]
does not imply
\[
P\neq NP.
\]
A true sentence may be unprovable in a particular system.  No truth--provability completeness
principle is assumed here.

 Evolutionary and classical complexity must not be conflated.
The separation in Theorem \ref{thm:evsep} concerns continuation-robust, history-sensitive classes.
A proof about standard Turing-machine \(P\) and \(NP\) is not automatically refuted by a
separation between different classes.  

 External distinguishability is compatible with internal indistinguishability.
An external metalanguage may distinguish predetermined from non-predetermined worlds by
quantifying over alternative histories.  The observer claim is only that inhabitants confined
to one realized history lack that access.





\section{Conclusion}

A mathematical object, similar to Brouwer's choice
sequences, may be well-defined because every realized value is unique and
persistent, without possessing one completed extension.
 

We then consider observers who are themselves part of the world they observe. All such observers interact with the same evolving world and therefore share a single realized history. Because they cannot return the entire world to exactly the same past state while retaining knowledge of what happened afterward, they cannot observe and compare two different continuations of that same past. Consequently, from the history available to them, the observers may be unable to determine whether the values they observe were fixed in advance or were progressively determined through interaction.

A persistently evolutionary oracle provides a concrete computational model of this idea. Within the history-sensitive semantics introduced in this paper, the model yields a separation between deterministic and nondeterministic polynomial-time computation. We then consider its consequence for the provability of \(P=NP\). Suppose that an observer cannot distinguish its actual world from an evolutionary world in which \(P\neq NP\), and suppose that the observer would carry out exactly the same reasoning and obtain the same purported proof of \(P=NP\) in both worlds. Since \(P=NP\) is false in the evolutionary world, that reasoning cannot constitute a sound proof of \(P=NP\).

 

The principal conclusion is therefore conditional and epistemic:
\[
  \boxed{\text{embedded observers cannot soundly prove \(P=NP\).}}
\]

\begin{quote}
    If we are embedded observers in a world satisfying the stated assumptions, this limitation applies to us regardless of whether the actual world is predetermined or non-predetermined. The argument does not require the actual world to be non-predetermined; it requires only that, from our position within the world, we cannot distinguish a predetermined world from an observationally compatible non-predetermined one.

\end{quote}

The argument does not prove classical \(P\neq NP\).  Its proposed contribution lies instead
in the foundations of computation and mathematical knowledge: it isolates predetermination as
an assumption distinct from well-definedness and asks what mathematical proof can establish
when the reasoner is itself a process inside one irreversible computational history.

%\section*{Acknowledgments} The author used generative AI tools for English-language editing and stylistic improvement. The mathematical ideas, definitions, arguments, and results presented in this paper are the author's own.
 

\begin{thebibliography}{99}

\bibitem{AaronsonWigderson2009}
S. Aaronson and A. Wigderson,
Algebrization: A new barrier in complexity theory,
\emph{ACM Transactions on Computation Theory} \textbf{1} (2009), no.~1, Article 2.

\bibitem{AroraBarak2009}
S. Arora and B. Barak,
\emph{Computational Complexity: A Modern Approach},
Cambridge University Press, Cambridge, 2009.

\bibitem{BakerGillSolovay1975}
T. Baker, J. Gill, and R. Solovay,
Relativizations of the \(P=?NP\) question,
\emph{SIAM Journal on Computing} \textbf{4} (1975), 431--442.

\bibitem{Bostrom2003}
N. Bostrom,
Are you living in a computer simulation?,
\emph{The Philosophical Quarterly} \textbf{53} (2003), 243--255.

\bibitem{Goldin2000}
D. Q. Goldin,
Persistent Turing machines as a model of interactive computation,
in \emph{Foundations of Information and Knowledge Systems},
Lecture Notes in Computer Science, vol.~1762, Springer, 2000, pp.~116--135.

\bibitem{GoldinSmolkaWegner2006}
D. Goldin, S. Smolka, and P. Wegner (eds.),
\emph{Interactive Computation: The New Paradigm},
Springer, Berlin, 2006.

\bibitem{GoldinWegner2008}
D. Goldin and P. Wegner,
The interactive nature of computing: Refuting the strong Church--Turing thesis,
\emph{Minds and Machines} \textbf{18} (2008), 17--38.

\bibitem{RazborovRudich1997}
A. A. Razborov and S. Rudich,
Natural proofs,
\emph{Journal of Computer and System Sciences} \textbf{55} (1997), 24--35.

\bibitem{TroelstraVanDalen1988}
A. S. Troelstra and D. van Dalen,
\emph{Constructivism in Mathematics: An Introduction}, Vols.~I--II,
North-Holland, Amsterdam, 1988.

\bibitem{VanAtten2007}
M. van Atten,
\emph{Brouwer Meets Husserl: On the Phenomenology of Choice Sequences},
Springer, Dordrecht, 2007.

\bibitem{WongEtAl2023}
M. L. Wong, C. E. Cleland, D. Arend Jr., S. Bartlett,
H. J. Cleaves II, H. Demarest, A. Prabhu, J. I. Lunine, and R. M. Hazen,
On the roles of function and selection in evolving systems,
\emph{Proceedings of the National Academy of Sciences}
\textbf{120} (2023), no.~43, e2310223120.

\end{thebibliography}


\appendix
\section{A Growing-Code Implementation of \texorpdfstring{\(\mathcal A\)}{A}}
\label{app:growing-code}

The process \(\mathcal A\) was defined in Section~\ref{sec:oracleA} by means of a
persistently growing finite automaton.  The same process can be represented
more concretely by a program whose code grows as new inputs are processed.
We describe here a simple  implementation.

The purpose of this implementation is not to introduce a different oracle.
Rather, it provides another representation of the same process
\(\mathcal A\).  States of the automaton are represented by directories,
transitions are represented by subdirectories, and the accepting status of
a state is represented by a small file stored in the corresponding
directory.

\subsubsection*{The directory structure}

The generated program is organized as a binary tree of directories.  Each
directory corresponds to one finite binary string, interpreted as a prefix
of an input.

The root directory, denoted by
\texttt{A\_code/}, corresponds to the empty string \(\epsilon\).  Its
subdirectory \texttt{0/}, when present, corresponds to the prefix \(0\),
and its subdirectory \texttt{1/} corresponds to the prefix \(1\).
Recursively, if a directory represents a prefix \(u\), then its
subdirectories
\[
\texttt{0/}
\qquad\text{and}\qquad
\texttt{1/}
\]
represent \(u0\) and \(u1\), respectively.

For example, if the path corresponding to the string \(101\) has already
been generated, the relevant part of the program has the form
\begin{verbatim}
A_code/
    node.py
    1/
        node.py
        0/
            node.py
            1/
                node.py
\end{verbatim}

Thus, the successive directories represent
\[
\epsilon,\qquad
1,\qquad
10,\qquad
101.
\]

\subsubsection*{The contents of \texttt{node.py}}

Every directory contains a small file called \texttt{node.py}.  The file
contains exactly three binary values:
\[
(A,Z,O)\in\{0,1\}^3.
\]

Their meanings are as follows:
\begin{enumerate}[label=(\roman*)]
    \item \(A=1\) if and only if the prefix represented by the current
    directory is accepting;

    \item \(Z=1\) if and only if the current directory has a
    \texttt{0/} subdirectory;

    \item \(O=1\) if and only if the current directory has a
    \texttt{1/} subdirectory.
\end{enumerate}

For example, a file may have the form
\begin{verbatim}
ACCEPTING = 0
ZERO_CHILD = 1
ONE_CHILD = 0
\end{verbatim}

This means that the corresponding prefix is currently nonaccepting, that
its \(0\)-extension has already been generated, and that its
\(1\)-extension has not yet been generated.

Hence each \texttt{node.py} has constant size.  The information accumulated
by \(\mathcal A\) is not stored in one growing table that must be searched.
Instead, it is distributed over the growing binary directory structure.

The correspondence with the automaton representation is therefore:
\[
\begin{array}{c|c}
\text{Automaton representation}
&
\text{Growing-code representation}
\\ \hline
\text{state corresponding to }u
&
\text{directory corresponding to }u
\\
u\xrightarrow{0}u0
&
\texttt{0/} subdirectory
\\
u\xrightarrow{1}u1
&
\texttt{1/} subdirectory
\\
u\in F
&
\texttt{ACCEPTING = 1}
\end{array}
\]

\subsubsection*{Processing an input}

Let
\[
x=a_1a_2\cdots a_n,
\qquad a_i\in\{0,1\}.
\]

The program begins in the root directory.  After reading \(a_1\), it
examines the corresponding bit in the root \texttt{node.py}.  If the
required subdirectory exists, the computation moves into that directory.
It then reads \(a_2\) and repeats the same procedure.

Thus, as long as all required branches already exist, the program visits
only the directories corresponding to
\[
\epsilon,\quad
a_1,\quad
a_1a_2,\quad
\ldots,\quad
a_1a_2\cdots a_n.
\]

In particular, the program never searches through directories belonging to
prefixes unrelated to the current input.

The behavior at the end of this traversal corresponds exactly to
Rules~1--3 of Section~5.

\subsubsection*{Implementation of Rule~1}

Suppose that the complete path corresponding to \(x\) already exists.
After reading the last symbol \(a_n\), the program reaches the directory
representing \(x\).

If its \texttt{node.py} contains
\begin{verbatim}
ACCEPTING = 1
\end{verbatim}
then the program returns \texttt{True}.  No directory or file is modified.

This is exactly Rule~1 of the automaton representation.

Because the accepting bit is never changed from \(1\) back to \(0\), the
answer remains permanent.

\subsubsection*{Implementation of Rule~2}

Suppose that the complete path corresponding to \(x\) exists, but the
terminal file contains
\begin{verbatim}
ACCEPTING = 0
\end{verbatim}

The program then needs to determine whether the state corresponding to
\(x\) has a one-symbol transition to an accepting state.

There are at most two possibilities: \(x0\) and \(x1\).

If
\texttt{ZERO\_CHILD = 1}, the program reads the
\texttt{node.py} contained in the \texttt{0/} subdirectory and checks its
accepting bit.  Similarly, if
\texttt{ONE\_CHILD = 1}, it checks the file in the \texttt{1/}
subdirectory.

If either immediate child has
\begin{verbatim}
ACCEPTING = 1
\end{verbatim}
then \(\mathcal A\) returns \texttt{False}.  The current file does not need
to be modified.

Otherwise, neither one-symbol extension is accepting.  In this case,
Rule~2 requires the current state to become accepting.  The program
therefore changes only one bit in the current \texttt{node.py}:
\begin{verbatim}
ACCEPTING = 0
\end{verbatim}
is replaced by
\begin{verbatim}
ACCEPTING = 1.
\end{verbatim}
The program then returns \texttt{True}.

Thus Rule~2 modifies at most one existing node file.

\subsubsection*{Implementation of Rule~3}

Suppose that while reading
\[
x=a_1\cdots a_n
\]
the program has already read
\[
a_1\cdots a_i
\]
but discovers that the subdirectory corresponding to \(a_{i+1}\) does not
exist.

This is the code representation of an undefined transition in the
automaton, and therefore Rule~3 applies.

The program creates the missing directory corresponding to
\(a_1\cdots a_{i+1}\), together with its \texttt{node.py}.  It then updates
the appropriate child bit in the parent file.

For example, if \(a_{i+1}=0\), the parent file is changed from
\begin{verbatim}
ZERO_CHILD = 0
\end{verbatim}
to
\begin{verbatim}
ZERO_CHILD = 1.
\end{verbatim}

The program repeats this construction for every remaining symbol of
\[
a_{i+1}\cdots a_n.
\]

The last generated directory represents the complete input \(x\).  Its
\texttt{node.py} is created with
\begin{verbatim}
ACCEPTING = 1
\end{verbatim}
and the program returns \texttt{True}.

Thus Rule~3 literally enlarges the generated program by adding the missing
branch of the binary directory tree.

\subsubsection*{Example}

Suppose that initially the program consists only of
\begin{verbatim}
A_code/
    node.py
\end{verbatim}
with
\begin{verbatim}
ACCEPTING = 0
ZERO_CHILD = 0
ONE_CHILD = 0
\end{verbatim}

and the first query is
\[
101.
\]

The required branch does not exist.  Rule~3 therefore creates
\begin{verbatim}
A_code/
    node.py
    1/
        node.py
        0/
            node.py
            1/
                node.py
\end{verbatim}

The final file, corresponding to \(101\), has
\begin{verbatim}
ACCEPTING = 1.
\end{verbatim}

The file corresponding to \(10\), however, has
\begin{verbatim}
ACCEPTING = 0
ZERO_CHILD = 0
ONE_CHILD = 1
\end{verbatim}

because it has a \(1\)-child leading to the accepting node \(101\).

Consequently, if \(10\) is queried next, the program reaches the directory
for \(10\), observes that it is nonaccepting, checks its immediate
\(1\)-child, finds that \(101\) is accepting, and returns
\texttt{False} by Rule~2.

By contrast, if \(10\) had been queried first in the initial state,
Rule~3 would have generated the path for \(10\) and created its final file
with
\begin{verbatim}
ACCEPTING = 1.
\end{verbatim}

Hence the growing-code representation reproduces the same history
dependence as the automaton representation.

\subsubsection*{Polynomial running time}

We finally discuss the running time of this implementation.

Let
\[
n=|x|.
\]

When processing \(x\), the program never scans the entire directory tree.
At every level it reads only the \texttt{node.py} corresponding to the
current prefix and then follows at most one of two possible subdirectories.
Therefore at most
\[
n+1
\]
nodes are visited while reading the input.

Each \texttt{node.py} contains only three binary values and therefore has
constant size.

If Rule~1 applies, the program performs only the traversal of the path of
\(x\), requiring \(O(n)\) node operations.

If Rule~2 applies, after this traversal the program examines at most two
additional node files, namely those corresponding to \(x0\) and \(x1\).
It may then modify one binary value in the file corresponding to \(x\).
Thus Rule~2 adds only a constant number of node operations.

If Rule~3 applies after \(i\) symbols have been read, at most
\[
n-i\le n
\]
new directories and constant-size \texttt{node.py} files are created.
Hence at most \(O(n)\) additional node operations are required.

Consequently, under the abstract cost model in which reading, creating, or
updating one constant-size node file is a unit-cost operation, the total
running time is
\[
T_{\mathcal A}(x)=O(n).
\]

For a literal filesystem implementation, accessing a directory at depth
\(i\)  requires processing a pathname of length \(O(i)\). Therefore, using the
conservative bound \(O(n)\) for each such access and noting that at most
\(O(n)\) nodes are accessed or created, we obtain
\(
T_{\mathcal A}(x)
\) is for sure polynomial time. 
Thus, under either cost model, the implementation runs in polynomial time
in the length of the current input.

Importantly, this bound is independent of the total number of nodes
generated during previous interactions.  The complete program may grow
without bound over time, but processing an input \(x\) requires access only
to the branch determined by the prefixes of \(x\).  Consequently, code
generated on unrelated branches does not increase the asymptotic running
time of the current query.

This growing-code representation therefore provides a concrete
polynomial-time implementation of the persistently evolutionary process
\(\mathcal A\).

\section{Proof of Theorem~\ref{thm:evsep}}
\label{app:evolutionary-separation}

In this appendix, we give the complete proof of
Theorem~\ref{thm:evsep}. The proof uses the fact that all computations
interact with the same persistent oracle state and that no fixed order is
imposed on admissible interactions with the oracle. Thus, while simulating
a candidate deterministic procedure, we may perform additional oracle
queries before allowing one of its queries to be processed. All such
queries become part of the same global history, and every answer already
returned by \(\mathcal A\) remains permanent.


Recall that
\[
L_{\mathcal A}
=
\{x\in\{0,1\}^* :
\mathcal A \text{ accepts }x\},
\]
and
\[
L_{\exists}
=
\left\{
w\in\{0,1\}^* :
\exists x\in\{0,1\}^*
\bigl(
|x|=|w|
\land
x\in L_{\mathcal A}
\bigr)
\right\}.
\]

We first establish two consequences of the evolutionary rules.

\begin{lemma}[Protecting a string]
\label{lem:protect-string}
Let \(H\) be the current history, and let
\(x\in\{0,1\}^n\). If no previously processed query has \(x\) as a prefix and the path corresponding to \(x0\) has not already been determined, then querying \(x0\) before \(x\) makes \(x0\) accepted and x rejected. 
Moreover, both answers remain unchanged in every continuation of \(H\).
\end{lemma}

\begin{proof}
The automaton does not yet contain a complete path labelled
\(x0\). Therefore, when \(x0\) is queried, Rule~3 adds the missing
part of the path and makes its final state accepting. Hence
\[
x0\in L_{\mathcal A}.
\]

Now consider a subsequent query of \(x\). The path labelled \(x\)
ends at a nonaccepting state having a transition labelled \(0\) to
the accepting state reached by \(x0\). Rule~2 therefore causes
\(\mathcal A\) to reject \(x\), so
\[
x\notin L_{\mathcal A}.
\]

By the persistence property of \(\mathcal A\), neither answer can
change in any later history.
\end{proof}

\begin{lemma}[Acceptance of an untouched string]
\label{lem:untouched-string}
Let \(H\) be a history, and let \(z\in\{0,1\}^n\). Suppose that no
query made in \(H\) has \(z\) as its length-\(n\) prefix. Then,
if \(z\) is queried next, \(\mathcal A\) accepts \(z\).
\end{lemma}

\begin{proof}
Since no previous query has \(z\) as its length-\(n\) prefix, no
previous interaction has created a complete path labelled \(z\) or
a path passing through \(z\) to one of its extensions. Consequently,
the automaton crashes at some undefined transition while reading
\(z\). Rule~3 then adds a path for the unread suffix and makes its
final state accepting. Thus,
\[
z\in L_{\mathcal A}.
\]
\end{proof}

\begin{theorem}[Evolutionary separation]
\label{thm:evsep-appendix}
Under continuation-robust semantics,
\[
P_{\mathrm{ev}}^{\mathcal A}
\neq
NP_{\mathrm{ev}}^{\mathcal A}.
\]
\end{theorem}

\begin{proof}
By Proposition~\ref{prop:belongtoNP},
\[
L_{\exists}\in NP_{\mathrm{ev}}^{\mathcal A}.
\]
Suppose, toward a contradiction, that
\[
L_{\exists}\in P_{\mathrm{ev}}^{\mathcal A}.
\]
Then there exists a deterministic oracle procedure \(D\) and a
polynomial \(p\) such that, for every input \(w\),

\begin{enumerate}[label=(\roman*)]
    \item \(D\) terminates within \(p(|w|)\) steps;
    \item \(D\)'s output is correct;
    \item \(D\)'s output is continuation-robust.
\end{enumerate}

Notethat all observers and computational procedures interact with the same global oracle state. Consequently, the computation of \(D\) need not occur in isolation: while \(D\) is running, another embedded observer may submit additional queries to \(\mathcal A\). These interactions modify the same persistent state subsequently encountered by \(D\). Hence the sequence consisting of \(D\)'s oracle queries together with the observer's additional queries forms one admissible global history.

\vspace{0.1cm}

Let \(H_0\) be the current global history, and let
\[
m
=
\max
\left\{
|u| :
u \text{ has previously been queried in }H_0
\right\},
\]
where \(m=0\) if no query has yet been made.

Choose \(n>m\) sufficiently large that
\[
p(n)<2^n.
\]
Fix an arbitrary string
\[
w\in\{0,1\}^n.
\]

We run \(D\) on \(w\) as one process within the   shared oracle \(\mathcal A\) and interleave its oracle interactions with queries submitted by another observer.

 

\begin{quote}
\textit{Whenever \(D\) is about to submit an oracle query \(y\) with
\(|y|\ge n,\)
let
\(x=\operatorname{pref}_n(y)\)
be the prefix of \(y\) of length \(n\). If \(x\) has not already been
protected, we, acting as another observer in the shared computational world,
first query \(x0\). The query \(y\) submitted by \(D\) is then processed
against the resulting global oracle state. We call such a string \(x\)
\emph{protected}.}
\end{quote}
By Lemma~\ref{lem:protect-string}, every protected string
\(x\in\{0,1\}^n\) is rejected by \(\mathcal A\), and this rejection is
permanent.

Since \(D\) executes for at most \(p(n)\) steps, it makes at most
\(p(n)\) oracle queries. Therefore, at most \(p(n)\) distinct strings
of length \(n\) become protected during the computation. Since
\(
p(n)<2^n,
\)
there exists at least one string
\[
z\in\{0,1\}^n
\]
that is not protected. Equivalently, no oracle query made during the
simulation of \(D\) has \(z\) as its length-\(n\) prefix.

After \(D\) terminates, there are two cases.

\medskip

\noindent
\textbf{Case 1: \(D\) rejects \(w\).}

Choose an unprotected string
\[
z\in\{0,1\}^n.
\]
Because no query made during the computation of \(D\) has \(z\) as
its length-\(n\) prefix, Lemma~\ref{lem:untouched-string} applies.
Querying \(z\) causes \(\mathcal A\) to accept it. Hence
\[
z\in L_{\mathcal A}.
\]

Since \(|z|=|w|=n\), it follows that
\[
w\in L_{\exists}.
\]

Thus, there is an admissible continuation of the history in which
the negative output of \(D\) is false. This contradicts the
continuation-robustness of \(D\)'s decision.

\medskip

\noindent
\textbf{Case 2: \(D\) accepts \(w\).}

Every protected string of length \(n\) has already been forced to be
rejected.

For each remaining unprotected string
\[
x\in\{0,1\}^n,
\]
query \(x0\). Since no previous oracle query has \(x\) as its
length-\(n\) prefix, Lemma~\ref{lem:protect-string} applies. Therefore
each such \(x\) is rejected.

After these finitely many additional interactions, every string of
length \(n\) is rejected:
\[
L_{\mathcal A}\cap\{0,1\}^n
=
\varnothing.
\]

Consequently, there exists no \(x\) satisfying
\[
|x|=|w|
\qquad\text{and}\qquad
x\in L_{\mathcal A}.
\]
Thus,
\[
w\notin L_{\exists}.
\]

Therefore, there is an admissible continuation of the history in which the positive output of \(D\) is false. This again contradicts continuation robustness.
Since \(D\) is assumed to decide \(L_{\exists}\) correctly in the shared computational environment, its correctness must be preserved under every admissible interleaving of interactions with the global oracle state. To refute the existence of such a decider, it is therefore sufficient to construct one admissible interleaving in which the output of \(D\) is not continuation-robust.



\medskip

In both cases, the output of \(D\) can be falsified by an admissible
continuation of the shared oracle history. Hence no deterministic
polynomial-time procedure continuation-robustly decides
\(L_{\exists}\). Therefore,
\[
L_{\exists}\notin P_{\mathrm{ev}}^{\mathcal A}.
\]

Since
\[
L_{\exists}\in NP_{\mathrm{ev}}^{\mathcal A},
\]
we conclude that
\[
P_{\mathrm{ev}}^{\mathcal A}
\neq
NP_{\mathrm{ev}}^{\mathcal A}.
\]
\end{proof}



\end{document}


 from __future__ import annotations

import importlib.util
from pathlib import Path
from uuid import uuid4


# ================================================================
# Persistent source-code directory
# ================================================================
#
# Everything that A learns is stored here AS PYTHON SOURCE CODE.
#
# Initially:
#
#     A_code/
#         node.py
#
# After learning, for example, the path 101:
#
#     A_code/
#         node.py
#         1/
#             node.py
#             0/
#                 node.py
#                 1/
#                     node.py
#
# Thus the SOURCE CODE itself grows as a binary tree.
#
# There is no persistent Python dictionary representing the tree.
# ================================================================

CODE_ROOT = Path("A_code")


# ================================================================
# Code of one node
# ================================================================
#
# Every prefix of a binary string corresponds to one tiny Python
# program.
#
# For example:
#
#       A_code/1/0/node.py
#
# represents the prefix "10".
#
# Each node contains:
#
#   1. whether this prefix is accepting;
#
#   2. a nested conditional telling us whether a 0-child or
#      a 1-child has already been generated.
#
# Importantly, each node.py remains CONSTANT SIZE.
# ================================================================


def node_source(
    accepting: bool = False,
    has_zero: bool = False,
    has_one: bool = False,
) -> str:
    """
    Generate the Python source code for one node of A.

    The generated source itself contains the conditional structure.
    """

    return f'''# Automatically generated part of process A.
#
# This file represents one prefix in A's growing source-code tree.
#
# It may be rewritten by A, but information is never removed.

ACCEPTING = {accepting}


def next_branch(bit):
    """
    Determine whether the requested branch has already been generated.
    """

    if bit == "0":
        return {has_zero}

    elif bit == "1":
        return {has_one}

    else:
        raise ValueError("bit must be 0 or 1")
'''


# ================================================================
# Initialization
# ================================================================

def initialize_A(reset: bool = False) -> None:
    """
    Initialize process A.

    If reset=True, delete the previously generated code and start
    again from the initial empty automaton.

    Normally reset=False should be used because A is persistent.
    """

    if reset and CODE_ROOT.exists():
        import shutil
        shutil.rmtree(CODE_ROOT)

    CODE_ROOT.mkdir(parents=True, exist_ok=True)

    root_file = CODE_ROOT / "node.py"

    if not root_file.exists():
        root_file.write_text(
            node_source(
                accepting=False,
                has_zero=False,
                has_one=False,
            ),
            encoding="utf-8",
        )


# ================================================================
# Loading ONE code node
# ================================================================

def load_node(node_directory: Path):
    """
    Execute one generated Python source file.

    Only this one small node.py file is loaded.

    We do NOT load the whole previously generated code tree.

    Consequently, processing a length-n input requires loading only
    the nodes lying on its branch.
    """

    node_file = node_directory / "node.py"

    module_name = "_A_node_" + uuid4().hex

    spec = importlib.util.spec_from_file_location(
        module_name,
        node_file,
    )

    module = importlib.util.module_from_spec(spec)

    assert spec.loader is not None
    spec.loader.exec_module(module)

    return module


# ================================================================
# Reading the current state of ONE code node
# ================================================================

def read_node(node_directory: Path):
    """
    Return:

        accepting,
        has_zero_child,
        has_one_child

    by EXECUTING the Python source code of this node.
    """

    module = load_node(node_directory)

    accepting = module.ACCEPTING

    has_zero = module.next_branch("0")
    has_one = module.next_branch("1")

    return accepting, has_zero, has_one


# ================================================================
# Rewriting ONE small piece of source code
# ================================================================

def rewrite_node(
    node_directory: Path,
    *,
    accepting: bool,
    has_zero: bool,
    has_one: bool,
) -> None:
    """
    Rewrite exactly one node of A's source code.

    Notice that this does NOT rewrite the whole code tree.

    The file being rewritten has constant size.
    """

    node_file = node_directory / "node.py"

    node_file.write_text(
        node_source(
            accepting=accepting,
            has_zero=has_zero,
            has_one=has_one,
        ),
        encoding="utf-8",
    )


# ================================================================
# Add one new branch to the SOURCE CODE
# ================================================================

def create_child(
    parent_directory: Path,
    bit: str,
) -> Path:
    """
    Grow A's source code by one level.

    Suppose parent_directory represents prefix u.

    If bit == "0", this function creates source code for prefix u0.

    If bit == "1", it creates source code for prefix u1.

    It also modifies the Python code of the parent so that its
    nested conditional structure now knows that this branch exists.
    """

    if bit not in {"0", "1"}:
        raise ValueError("bit must be 0 or 1")

    accepting, has_zero, has_one = read_node(parent_directory)

    child_directory = parent_directory / bit
    child_directory.mkdir(parents=True, exist_ok=True)

    child_file = child_directory / "node.py"

    # Create the new source-code node.
    if not child_file.exists():
        child_file.write_text(
            node_source(
                accepting=False,
                has_zero=False,
                has_one=False,
            ),
            encoding="utf-8",
        )

    # ------------------------------------------------------------
    # Now MODIFY THE PARENT'S SOURCE CODE.
    #
    # Before:
    #
    #     if bit == "0":
    #         return False
    #
    # After learning branch 0:
    #
    #     if bit == "0":
    #         return True
    #
    # Thus the program itself has grown.
    # ------------------------------------------------------------

    if bit == "0":
        has_zero = True
    else:
        has_one = True

    rewrite_node(
        parent_directory,
        accepting=accepting,
        has_zero=has_zero,
        has_one=has_one,
    )

    return child_directory


# ================================================================
# Make one prefix accepting
# ================================================================

def make_accepting(node_directory: Path) -> None:
    """
    Permanently modify the source code so that this prefix becomes
    accepting.

    Accepting status is never changed back to False.
    """

    accepting, has_zero, has_one = read_node(node_directory)

    if accepting:
        return

    rewrite_node(
        node_directory,
        accepting=True,
        has_zero=has_zero,
        has_one=has_one,
    )


# ================================================================
# Check whether a node has an ACCEPTING one-symbol extension
# ================================================================

def has_accepting_child(node_directory: Path) -> bool:
    """
    Implements the test required by Rule 2.

    We inspect at most TWO child programs.

    Therefore this operation has constant cost with respect to the
    total size of A's accumulated code.
    """

    _, has_zero, has_one = read_node(node_directory)

    if has_zero:
        child0 = node_directory / "0"
        accepting0, _, _ = read_node(child0)

        if accepting0:
            return True

    if has_one:
        child1 = node_directory / "1"
        accepting1, _, _ = read_node(child1)

        if accepting1:
            return True

    return False


# ================================================================
# RULE 3
# ================================================================

def grow_missing_suffix(
    current_directory: Path,
    suffix: str,
) -> Path:
    """
    Rule 3.

    The current code cannot read the remaining suffix.

    We therefore physically GENERATE new Python source code for
    every remaining bit.

    If suffix has length k, at most k new constant-size Python
    programs are generated.

    Hence growth caused by this query is O(k), and therefore O(|x|)
    in the abstract source-update model.
    """

    current = current_directory

    for bit in suffix:
        current = create_child(current, bit)

    # Rule 3:
    # the final newly generated prefix becomes accepting.
    make_accepting(current)

    return current


# ================================================================
# PROCESS A
# ================================================================

def A(x: str) -> bool:
    """
    Persistent evolutionary process A.

    Input
    -----
    x = a1 a2 ... an, where ai is 0 or 1.

    Output
    ------
    True  = YES / ACCEPT
    False = NO  / REJECT


    ==============================================================
    IMPORTANT
    ==============================================================

    A does NOT search through everything it has previously learned.

    It follows only the branch

          a1 -> a2 -> ... -> an.

    Thus an input of length n visits at most n+1 code nodes.

    The total amount of source code generated during the entire
    history may be extremely large, but the computation on x does
    not inspect unrelated branches.
    """

    if not isinstance(x, str):
        raise TypeError("x must be a string")

    if any(bit not in {"0", "1"} for bit in x):
        raise ValueError("x must be a binary string")

    initialize_A()

    # Start at the program representing the empty prefix.
    current = CODE_ROOT

    # ============================================================
    # Read x symbol by symbol.
    # ============================================================

    for i, bit in enumerate(x):

        node = load_node(current)

        # --------------------------------------------------------
        # This is the nested conditional structure belonging to
        # the CURRENT prefix.
        #
        # Does the code already contain the required branch?
        # --------------------------------------------------------

        branch_exists = node.next_branch(bit)

        if branch_exists:

            # The code already knows this branch.
            # Follow it.
            current = current / bit

        else:

            # ====================================================
            # RULE 3
            # ====================================================
            #
            # The current source code cannot finish reading x.
            #
            # The unread suffix is:
            #
            #       x[i:]
            #
            # A now MODIFIES ITS OWN SOURCE CODE and generates
            # exactly the missing nested branch.
            # ====================================================

            grow_missing_suffix(
                current_directory=current,
                suffix=x[i:],
            )

            return True

    # ============================================================
    # We have successfully read all of x.
    #
    # Therefore Rule 1 or Rule 2 applies.
    # ============================================================

    accepting, _, _ = read_node(current)

    # ============================================================
    # RULE 1
    # ============================================================

    if accepting:

        # x was already accepted.
        #
        # Nothing changes.
        return True

    # ============================================================
    # RULE 2 -- rejection case
    # ============================================================

    if has_accepting_child(current):

        # x is nonaccepting and one of x0 or x1 is accepting.
        #
        # Therefore x is rejected.
        #
        # Notice that we do not need to store a separate rejection
        # flag: the persistent accepting child permanently guarantees
        # that Rule 2 will reject x again in every future query.
        return False

    # ============================================================
    # RULE 2 -- acceptance case
    # ============================================================

    # x is nonaccepting and has no accepting one-symbol extension.
    #
    # Modify x's source code and make it accepting permanently.
    make_accepting(current)

    return True









    \begin{proposition}[Polynomial-time response of the growing-code implementation]
\label{prop:growing-code-polytime}
Let \(x\in\{0,1\}^*\) with \(n=|x|\). Assume that reading, writing,
or loading a constant-size source file located at a path of length \(O(n)\)
takes \(O(n)\) time. Then the growing-code implementation of
\(\mathcal A\) processes the query \(x\) in \(O(n^2)\) time.

In particular, the running time of a query is polynomial in the length of
the current input and does not depend on the total number of code nodes
generated during previous interactions.
\end{proposition}

\begin{proof}
Every node of the growing program corresponds to a prefix of some previously
processed input.  The source file associated with each node has constant
size: it contains only an accepting flag and two conditional branches,
corresponding to symbols \(0\) and \(1\).  Hence the size of an individual
node does not increase as the global program evolves.

Consider a query
\[
x=a_1a_2\cdots a_n.
\]
The computation begins at the root and follows the unique sequence of
prefixes
\[
\epsilon,\quad
a_1,\quad
a_1a_2,\quad
\ldots,\quad
a_1a_2\cdots a_n.
\]
Consequently, at most \(n+1\) existing code nodes are visited.

At a node corresponding to a prefix of length \(i\le n\), the program
performs only a constant number of operations: it loads the corresponding
constant-size source file and tests whether the branch labelled by the next
input symbol exists.  The path identifying that source file has length
\(O(i)\), and therefore at most \(O(n)\).  Under the stated cost assumption,
processing one visited node consequently requires \(O(n)\) time.  Since at
most \(n+1\) nodes are visited, traversal of an already existing path costs
\[
O(n)\cdot O(n)=O(n^2).
\]

It remains to consider modifications of the source code.

If Rule~1 applies, no modification occurs.

If Rule~2 applies, the program inspects at most two immediate children of
the terminal node.  Each such operation concerns a constant-size source file
at a path of length \(O(n)\).  If the terminal node must be made accepting,
only one constant-size source file is rewritten.  Thus Rule~2 requires
\(O(n)\) additional time.

Suppose instead that Rule~3 applies after the first \(i\) symbols have been
read.  The unread suffix has length at most \(n-i\le n\).  The program
generates at most one new code node for each symbol of this suffix.
Therefore at most \(n\) new constant-size source files are created.  Each
creation or modification concerns a path whose length is \(O(n)\), and
hence costs \(O(n)\) under the stated assumption.  The total cost of growing
the missing branch is therefore at most
\[
O(n)\cdot O(n)=O(n^2).
\]

Combining traversal and possible modification gives
\[
T_{\mathcal A}(x)=O(n^2).
\]

Most importantly, no step requires traversing or searching the complete
source-code tree generated by earlier computations.  Branches corresponding
to prefixes unrelated to \(x\) are never inspected.  Hence the bound depends
only on \(n=|x|\), and not on the number of strings queried before \(x\).
Therefore every query is processed in polynomial time.
\end{proof}
